\documentclass[11pt]{article}
\usepackage{microtype,amsmath,amssymb,amsthm,mathtools,booktabs,array}
\usepackage[margin=1.05in]{geometry}
\usepackage[colorlinks=true,linkcolor=blue!45!black,citecolor=blue!45!black,urlcolor=blue!45!black]{hyperref}
\usepackage{xcolor}
\hypersetup{pdftitle={The dimension-degree boundary for Gaussian moment vanishing},pdfsubject={Working mathematical manuscript}}
\newtheorem{theorem}{Theorem}[section]
\newtheorem{proposition}[theorem]{Proposition}
\newtheorem{lemma}[theorem]{Lemma}
\newtheorem{corollary}[theorem]{Corollary}
\theoremstyle{definition}\newtheorem{definition}[theorem]{Definition}
\theoremstyle{remark}\newtheorem{remark}[theorem]{Remark}
\newcommand{\E}{\mathbb E}
\newcommand{\C}{\mathbb C}
\newcommand{\tr}{\operatorname{tr}}
\newcommand{\wt}{\operatorname{wt}}
\newcommand{\G}{\mathsf G}
\DeclareMathOperator{\Span}{span}
\title{Bounded critical models and the regular Gaussian seed}
\author{}\date{17 November 2026: internal written proof}
\begin{document}\maketitle
This note proves the regular-part reduction for polynomial Gaussian
phases and the complete nonisotropic zero-$F_2$ cusp,tangent and
concurrent-line normal forms. The boundary manuscript supplies the
separate nonzero-$F_2$ and isotropic-frame gates. The new argument
constructs no unnecessary escape characters:it proves their exclusion
from the regular summand,then fills that summand exactly.
Classical formal decomposition and direct-image rank mechanisms are
credited to Sabbah and Roucairol,and tame Milnor rank to Pham.
The geometric rank application uses complex polynomial coefficients;
no arbitrary nonconvergent family or all-ring extension is asserted.
No novelty,independent replay,formal verification or external review
is claimed.

Let $q$ be the full nondegenerate affine Gaussian after an invertible
affine change,with normalized formal expectation $\mathbb E$. Set
\[
 \Phi_A(t)=\sum_{j\ge0}\mathbb E[AP^j]t^j/j!,\qquad
 M=\mathbb C(t)[x]/\sum_i(\partial_i+tP_i-q_i)\mathbb C(t)[x].
\]
The connection is $\nabla_t=\partial_t+P$ and Gaussian integration by
parts gives a horizontal functional $J:M\to\mathbb C((t))$.
No convergence of the moment series,actual-period choice or image-kernel
computation is assumed.

\begin{lemma}[Rational seed restriction]\label{lem:seedfunctional}
If $J(a)=1$,the differential submodule $U$ generated by $a$ has a
horizontal quotient to $\mathbb C(t)$ taking $a$ to one.
\end{lemma}
\begin{proof}
The generator values are $1,0,0,\ldots$. Applying $J$ to every rational
relation between these generators shows that those values define a
well-defined horizontal rational functional. Its value at $a$ is
surjective. Formal extension at infinity is exact.
\end{proof}
\begin{lemma}[Local isolated characters]\label{lem:isolatedcharacters}
For an exact local phase $\alpha+sz+y^{\mu+1}$,an analytic amplitude,
and an analytic volume unit,there are $\mu$ horizontal formal
characters. With $k_0=\mu+1$,the character indexed $j$, $1\le j<k_0$,
vanishes on a monomial $y^a s^b z^c$ unless $b=c$ and
$a=j-1+\ell k_0$, $\ell\ge0$. Otherwise it is
\begin{equation}\label{eq:ordinarycharacters}
 e^{\alpha t}t^{-1-j/k_0}
       (-1)^{b+\ell}b!\,(j/k_0)_\ell\,t^{-b-\ell}.
\end{equation}
\end{lemma}
\begin{proof}
At any given order only finitely many Taylor coefficients contribute.
The recurrences $I_{a+k_0}=-(a+1)I_a/(k_0t)$ and
$J_{b+1,b+1}=-(b+1)J_{b,b}/t$ annihilate all spatial exact forms.
The same recurrences verify that time differentiation inserts the
phase. The leading values are the jets $1,y,\ldots,y^{\mu-1}$,
multiplied by the analytic amplitude and volume unit. Their normalized
jet matrix is triangular with nonzero diagonal. Thus all multiplicities
are retained,without a Morse or distinct-critical-value hypothesis.
\end{proof}
\subsection{It is enough to complete the regular part}
\label{sec:boundedregular}

The cusp and tangent cases share a substantial simplification.
We need not construct every escaping Gaussian character. It is
enough to exclude escaping contributions to the regular part,
compute that part's full rank, and construct precisely that many
local functionals. This keeps the original marker while avoiding
an unnecessary classification of the other exponential blocks.

\begin{lemma}[A positive seed on the regular part]
\label{lem:regularpositiveseed}
Let $M$ and a seed $a$ satisfy the formal Gaussian hypotheses of
Lemma~\ref{lem:seedfunctional}. After a finite cover at infinity,
write $M_0$ for the zero-exponential summand. Evaluate its dual
functionals on the canonical projection of $a$ to $M_0$. Suppose $M_0$ has a
full horizontal dual frame over $\C((\rho))$, and every member
takes $a$ into $\rho\C[[\rho]]$. Then the formal Gaussian
functional cannot take $a$ to one. No semisimplicity assumption
on the nonzero-exponential summands is required.

The same conclusion holds if the regular frame exists over
$\C((\rho))(\log\rho)$ and its seed values are finite
polynomials in $\log\rho$ with coefficients of strictly positive
$\rho$ order.
\end{lemma}
\begin{proof}
The rational seed restriction supplies a quotient $U\to\C(t)$.
The formal exponential decomposition is functorial and exact on
submodules and quotients; different exponential types have no
horizontal morphisms. These are the classical formal decomposition
facts, recalled in Sabbah~\cite[Section~3, Corollary~3.3 and its
proof]{SabbahStationary}. Consequently the quotient can only use
$U_0\subset M_0$. Since $M_0$ is trivial over the displayed field,
it is semisimple there, and this quotient extends to a horizontal
functional on $M_0$. It is a constant linear combination of its
full dual frame. Positive seed orders cannot sum to one.

Adjoining $\log\rho$, with derivative $1/\rho$, does not add
constants or negative-power exponential solutions. The same
argument applies over this differential field. A finite sum of
positive powers times logarithmic polynomials still cannot be
one. This also explains why a full regular frame, rather than
semisimplicity of the whole ambient module, is the relevant
hypothesis.
For an elementary check, expand this field in
$\C(\log\rho)((\rho))$. The operator $\rho\partial_\rho$
preserves its valuation. A nonzero exponential twist has a
negative-valuation scalar term in that operator, which cannot
be cancelled by a regular matrix. Its horizontal maps to the
trivial module therefore vanish even over this field. Its
constants are $\C$, since a rational function $a(z)$ satisfying
$a'+na=0$ is zero for integer $n\ne0$ and constant for $n=0$.
\end{proof}

\begin{proposition}[Bounded critical branches determine the regular rank]
\label{prop:boundedregularrank}
Let $P$ have degree $D\ge2$ and leading form $F_D$, and let $q$
be quadratic. Suppose, locally uniformly for small nonzero
$\epsilon$, the phase $S=P/\epsilon-q$ has a complex linear
gradient lower bound outside a compact set. The constants and
compact set may depend on $\epsilon$. Assume also
\begin{equation}\label{eq:nonisoleading}
 q_2(w)\ne0\quad\hbox{whenever }[w]\in\mathbb P^{n-1}
                           \hbox{ and }\nabla F_D(w)=0.
\end{equation}
Then the rank of the full zero-exponential summand equals the sum
of the generic local Milnor numbers of bounded phase-critical
branches whose limiting $P$ value is zero. An escaping branch
never has action polar part $\alpha/\epsilon$ for any constant
$\alpha$.
\end{proposition}
\begin{proof}
At a stationary branch the chain rule gives
$S'=-P/\epsilon^2$, hence $q=-S-\epsilon S'$.
A bounded branch has bounded $q$, so its action has only the
possible polar term $P(x(0))/\epsilon$. If a branch escapes,
divide $\nabla P=\epsilon\nabla q$ by $\|x\|^{D-1}$.
Its limiting direction is singular for $F_D$.
Condition~\eqref{eq:nonisoleading} makes the quadratic term of
$q$ along that direction nonzero. Thus $q(x(\epsilon))$ has a
pole, whereas $-S-\epsilon S'$ would be bounded for either zero
polar part or a single term $\alpha/\epsilon$. This excludes
both possibilities, without an escape-scale or Morse assumption.

We spell out the global rank comparison. Fixed-parameter tameness
gives full twisted rank equal to the total critical Milnor number
by Pham's decomposition. For the regular algebraic input
$\mathcal O$ on $\C^*_{\epsilon}\times\C^n$, form the ordinary
direct image by $(\epsilon,S)$. Its relevant cohomology is the
middle reduced Gauss--Manin module; the constant fiber cohomology
has zero exponential integral in the affine $S$ coordinate.
There is no additional finite-target discriminant from spatial
infinity. Indeed, on $S=g$ in a bounded target set,
\[
 \partial_\epsilon S=-(g+q)/\epsilon=O(1+\|x\|^2).
\]
The locally uniform gradient estimate gives a parameter-lifting
field of at most linear growth. The target-lifting field is
$O(1/\|x\|)$. Both are complete outside a compact set on each
parameter neighborhood. Thus their outer fiber pairs trivialize.
Every nonvertical discriminant curve at finite target comes
from an affine critical value, with conormal multiplicity equal
to the sum of its local Milnor numbers.

First ramify all critical-value curves into graphs. Roucairol's
nonramified rank formula, Theorem~0.1(2), Section~1.2, together
with Theorem~0.3~\cite{RoucairolFormal}, computes every nonzero
exponential-block rank from those graph multiplicities. Count
all conjugate graphs even when their polar parts agree. No
distinct-action assumption or monodromy formula is used.
Subtract these ranks from the full critical Milnor number.
Exactly the bounded-action branches remain. By the first part
they are precisely the bounded branches with limiting $P$ value
zero. This proves the assertion about the entire regular rank.
\end{proof}

We use the following explicit way to construct that regular part.
Suppose a bounded group at $P$ value $\alpha$ has coordinates,
after a cover, in which
\begin{equation}\label{eq:boundedfixedmodel}
 S=\alpha t+H_0(z)+E_\rho(z),\qquad
 \omega=\rho^\kappa J_\rho(z)\,dz,
 \quad E_\rho\in\rho\C[z][[\rho]],\quad\kappa>0,
\end{equation}
where $J_\rho$ is a unit with polynomial coefficients at each
order. If $\lambda$ is a dual functional on the fixed twisted
quotient of $H_0$, then
\[
 L(A)=e^{\alpha t}\rho^\kappa\lambda\!\left(
           e^{E_\rho}J_\rho A(x(\rho,z))\right)
\]
is a formal functional on the original module. The product rule
proves spatial exactness. Cartan's formula for the coordinate
velocity, together with the displayed moving Jacobian, proves
parameter horizontality. Every coefficient uses finitely many
polynomial terms. Fixed original markers with nonnegative powers
in these coordinates have positive seed order.

For completeness one must also prove that all bounded critical
branches in the group have bounded $z$, and that the fixed
critical scheme has total Milnor number equal to its twisted
rank. Conservation of local intersection multiplicity then
counts the whole group. A polynomial basis of the fixed quotient
gives independent functionals at leading order. Distinct groups
are separated by polynomial jets, with arbitrarily high vanishing
orders to overcome the finite negative powers introduced by the
scaled representatives. When their number equals the regular
rank of Proposition~\ref{prop:boundedregularrank}, they exhaust
$M_0$. These checks will be made explicitly below; the formal
perturbation formula alone does not establish completeness.

\subsection{The cusp and tangent families together}
\label{sec:cusptangent}

For the nonisotropic zero-$F_2$ alternatives retain the full Gaussian
\[
 q=\tfrac12(Au^2+Bv^2+mT^2+2Cuv+2DuT+2EvT)
                  +\ell_u u+\ell_v v+\ell_TT+q_0,
 \qquad m\ne0.
\]
Translations and linear changes give the two exact normal forms
\begin{align}
 P_C&=T(u^2+\gamma v+h)+v^3+buv+eu+fv+k,
                       &&\hbox{cusp},\label{eq:cuspnormal}\\
 P_T&=T(u^2+gv+h)+(u+c)v^2+eu+fv+k,
                       &&\hbox{tangent line--conic}.\label{eq:tangentnormal}
\end{align}
Here is the leading-form reduction. At the unique singular
point, $F_3$ is $T$ times a rank-one binary quadratic plus a
binary cubic. Choose that quadratic as $u^2$. A linear shift
of $T$ removes the $u^3,u^2v$ terms, leaving
$Tu^2+a uv^2+b_0v^3$. If $b_0\ne0$, a shift of $v$ by a
multiple of $u$ removes $uv^2$, and rescaling gives the cusp.
If $b_0=0$, tangency to a smooth conic requires $a\ne0$;
rescaling gives the tangent form. The direction of $T$ is
unchanged, so its Gaussian square coefficient remains nonzero.
For the cusp, translate $u$ to remove $uT$, translate $T$ to
remove $u^2$, and translate $v$ to remove $v^2$. For the tangent
case use the same first two translations, then a $v$ translation
to remove $uv$. These operations leave the displayed arbitrary
coefficients and transform, rather than reset, the means and
cross terms of $q$.

\begin{lemma}\label{lem:cusptangentgate}
For both leading forms, $F_2$ nonzero at the singular direction
implies $\|\nabla P\|\ge c\|x\|$ outside a complex compact set.
In the nonisotropic forms above, each fixed nonzero phase
$tP-q$ has the same type of lower bound, locally uniformly in
$t$. Their full generic twisted ranks are six and five,
respectively.
\end{lemma}
\begin{proof}
Any failure of a linear gradient bound approaches the unique
leading singular direction, so $u,v=o(T)$. For the cusp,
the $u$ equation divided by $T$ makes $u$ bounded. The $T$
equation divided by $T$ then has the nonzero limit $2F_2(0,0,1)$.
For the tangent, the transverse equations give
$uv=O(T)$ and $2Tu+v^2=O(T)$. If $u$ is unbounded these imply
$u=O(|T|^{1/3})$; otherwise $u^2/T\to0$ directly. The $T$
equation again has the same nonzero limit. Replacing the quadratic value by $-m/2$
gives the fixed-phase bounds for the normal forms. Compactness
of the coefficient estimates gives local parameter uniformity.

For clarity, the ranks do not come from an escape count. The
three quadratic gradient equations have total projective length
eight. Their only infinity point is the $T$ direction. In the
chart $T=1$ with homogenizing coordinate $z$, the cusp $u$
equation solves $u=O(z)$; the $T$ equation then gives $z=0$,
and the last equation is $3tv^2=0$. Its length is two. For
the tangent the $T$ equation solves $z=O(u^2)$, the $u$ equation
gives $u=-v^2/2+O(v^4)$, and the $v$ equation has leading term
$-tv^3$. Its length is three. Fixed-phase coercivity excludes
affine components of positive dimension. B\'ezout and Pham's
tame comparison therefore give ranks $8-2=6$ and $8-3=5$.
\end{proof}

Every escape in these two normal forms approaches the $T$
direction, where $q_2=m/2\ne0$. Proposition~\ref{prop:boundedregularrank}
therefore applies. We now list all possible bounded contributions;
the other exponential blocks require no explicit construction.

\subsubsection*{Isolated critical points on a smooth constraint}

Write $P=T A(u,v)+G(u,v)$. Where $A=0$ is smooth, take $s=A$
and a coordinate $y$ along that curve. The analytic shear
\[
 T'=T+\frac{G(y,s)-G(y,0)}s
\]
gives the exact identity $P=sT'+G(y,0)$. At an isolated critical
point this becomes $\alpha+sT'+y^{\mu+1}$. The already explicit
characters \eqref{eq:ordinarycharacters} provide all $\mu$
dimensions with positive original-marker order.

For the cusp, $\gamma\ne0$ gives a degree-six restriction to
$v=-(u^2+h)/\gamma$, so total isolated multiplicity is five.
If $\gamma=0,h\ne0$, the two lines $u=\pm\sqrt{-h}$ each
have a cubic restriction in $v$, giving multiplicity two each.
For the tangent, $g\ne0$ gives a degree-five restriction and
multiplicity four. If $g=0,h\ne0$, each line $u=u_i$ gives
\[
 (u_i+c)v^2+fv+eu_i+k.
\]
It has one Morse point when $u_i+c\ne0$, none when
$u_i+c=0,f\ne0$, and a critical parabola when both vanish.
The last case is included below. All repeated roots and equal
critical values in the isolated cases retain their Milnor
multiplicities; no Morse perturbation is assumed.

\subsubsection*{Four fixed models along the $T$ lines}

Here and below a scalar constant is harmless but retained in the
phase value. Set $\epsilon=1/t$. The critical equation in $T$
is
\begin{equation}\label{eq:cuspcriticalbound}
 u^2=\epsilon(mT+Du+Ev+\ell_T)
\end{equation}
whenever the constraint is the double line $u^2=0$.
Thus every bounded phase-critical branch has $u=O(\epsilon^{1/2})$.
This estimate will also verify that the following scalings lose
no bounded contribution.

First consider a critical $T$ line at $u=0,v=v_0$ whose transverse
quadratic part is $Tu^2+\lambda(v-v_0)^2+\beta u(v-v_0)$,
with $\lambda\ne0$. Scale $u=t^{-1/2}U$,
$v=v_0+t^{-1/2}V$, leaving $T$ fixed. The volume is $t^{-1}$
and the fixed phase is
\[
 H_0=TU^2+\lambda V^2+\beta UV-mT^2/2-\eta T-q_*.
\]
Complete $V$, then $T$. This is a quartic in $U$, with leading
coefficient $1/(2m)$, plus two nondegenerate Gaussians. Its
twisted rank and critical Milnor number are both three, including
all quartic collisions. Equation~\eqref{eq:cuspcriticalbound}
and the $v$ equation with nonzero linear coefficient $2\lambda$
give $v-v_0=O(t^{-1/2})$. Every bounded branch in this group has
bounded scaled variables. The model therefore supplies all three
local dimensions, with original-marker values in
$e^{\alpha t}t^{-1}\C[[t^{-1/2}]]$.

If instead $\lambda=0,\beta\ne0$, the same scaling gives
\[
 H_0=TU^2+\beta UV-mT^2/2-\eta T-q_*.
\]
The polynomial shear $V\mapsto V+TU/\beta$ removes $TU^2$.
This is a nondegenerate three-variable Gaussian, of rank one.
The $u$ equation with coefficient $\beta$ bounds $v-v_0$ by
$O(u)+O(\epsilon)$, so bounded branches are again captured.
It gives one dimension with the same positive volume order.

For the cusp's remaining line, $P=T u^2+v^3+k$, use
$u=t^{-1/2}U$, $v=t^{-1/3}V$. The fixed phase is
\[
 TU^2+V^3-mT^2/2-\ell_TT-q_0.
\]
Completing $T$ gives a quartic, a cubic and one Gaussian, of
total rank and Milnor number $3\cdot2=6$. The volume is
$t^{-5/6}$. Bounded critical branches satisfy $u=O(t^{-1/2})$
and, from $3v^2=\epsilon q_v$, $v=O(t^{-1/2})$, so their
scaled coordinates are bounded.

For the tangent's remaining line, $P=T u^2+uv^2+k$, use
$u=t^{-1/2}U$, $v=t^{-1/4}V$. Completing $T$ gives the fixed
two-variable phase
\[
 U^4/(2m)+UV^2-\ell_TU^2/m
\]
and one Gaussian. Its rank and critical Milnor number are five.
Here is a direct check: give $U,V$ weights $1/4,3/8$. The
leading derivatives are $2U^3/m+V^2$ and $2UV$, with sole
common zero the origin and quotient basis
$1,U,U^2,V,V^2$. They form a regular sequence. The associated
graded twisted de Rham complex is their Koszul complex,
concentrated in top cohomology; all derivative and lower-phase
terms have smaller weight. It follows that the same five
classes give the twisted rank. The ordinary gradient quotient
has the same weighted leading calculation, so its total
Milnor number is also five for every $\ell_T$. The volume is
$t^{-3/4}$. Equation~\eqref{eq:cuspcriticalbound} bounds $u$,
and $2Tu+v^2=\epsilon q_u$ gives $v=O(t^{-1/4})$. This
captures every bounded branch.

All remaining terms in these four scaled phases have positive
order in a common finite cover. Their coefficients are
polynomials at each order, and every fixed original marker has
nonnegative order. Formula~\eqref{eq:boundedfixedmodel} constructs
the stated complete local functionals. The Jacobian factors and
coordinate-velocity terms are retained, exactly as in its proof.

\subsubsection*{The critical parabola}

This occurs only for the tangent form, with $g=0$, $h\ne0$,
$c^2+h=0$, $f=0$, $c\ne0$. Put
\[
 x=u+c,\qquad s=(x-2c)T+v^2+e,\qquad\alpha=k-ec.
\]
Exactly $P=\alpha+xs$. Write $x=\rho X$, $s=\rho Y$,
$\rho=t^{-1/2}$. The volume is
\[
 \frac{\rho^2}{\rho X-2c}\,dX\,dv\,dY,
\]
and the fixed phase is $XY+\psi(v)$, where
\[
 \psi(v)=-q\!\left(-c,v,\frac{v^2+e}{2c}\right).
\]
Its quartic coefficient is $-m/(8c^2)\ne0$, so its fixed
twisted rank and critical Milnor number are three. At every
bounded phase-critical branch near this parabola, the exact
normal equations give $x=\epsilon q_s=O(\epsilon)$ and
$s=\epsilon q_x=O(\epsilon)$; the denominators stay away
from zero. Thus $X,Y$ are bounded. The scaled Jacobian is a
unit times $\rho^2$, and its expansion and that of the phase
have polynomial coefficients at each order. This gives all
three dimensions, with strictly positive original-marker order.
The other constraint line has its one isolated Morse point.

\subsubsection*{Exhaustive routing and completeness}

The remaining double-line cases require only the following two
tables. Their rows are disjoint. A missing common root means
there is no bounded contribution at all, not an omitted escape.

\begin{center}\small
\begin{tabular}{@{}p{0.36\textwidth}p{0.54\textwidth}@{}}
\toprule
Cusp with $\gamma=h=0$ & Bounded $P$-critical groups\\
\midrule
$f\ne0$ & For each root $3v_i^2+f=0$ with $bv_i+e=0$,
one line with $\lambda=3v_i$, $\beta=b$, rank three\\
$f=0$, $e\ne0$ & None, for every $b$\\
$f=e=0$, $b\ne0$ & One line with $\lambda=0$, $\beta=b$, rank one\\
$f=e=b=0$ & Quartic--cubic model, rank six\\
\bottomrule
\end{tabular}
\end{center}
The line values in the first row are
$\alpha_i=v_i^3+fv_i+k$; the other line values are $k$.

\begin{center}\small
\begin{tabular}{@{}p{0.36\textwidth}p{0.54\textwidth}@{}}
\toprule
Tangent with $g=h=0$ & Bounded $P$-critical groups\\
\midrule
$c\ne0$ & Put $v_0=-f/(2c)$. None if $e+v_0^2\ne0$;
otherwise one rank-three line with $\lambda=c$, $\beta=2v_0$,
value $k-f^2/(4c)$\\
$c=0$, $f\ne0$ & None, for every $e$\\
$c=f=0$, $e\ne0$ & Two lines $v_i^2=-e$, each with
$\lambda=0$, $\beta=2v_i$, rank one, value $k$\\
$c=f=e=0$ & Weighted quartic model, rank five, value $k$\\
\bottomrule
\end{tabular}
\end{center}

To verify these tables directly, the double constraint forces
$u=0$. The cusp equations reduce to $3v^2+f=0$, $bv+e=0$;
the tangent equations reduce to $2cv+f=0$, $v^2+e=0$.
The remaining $T$ coordinate is free. Together with the smooth
constraint and separated-line alternatives already proved,
this exhausts all coefficients of both normal forms.

For every fixed model the critical scheme is finite with the
stated total Milnor number. The bounds above show that every
bounded phase-critical branch has bounded scaled coordinates.
Conservation of local critical multiplicity thus gives exactly
the stated group counts. Proposition~\ref{prop:boundedregularrank}
shows that the groups at value zero account for the entire
regular rank; all escaping branches are excluded from it.

We also need independence on the original polynomial module.
Within a fixed group evaluate a polynomial basis of its fixed
quotient: its normalized leading matrix is invertible by
\eqref{eq:boundedfixedmodel}. Isolated groups use their finite
power jets. Different line groups have distinct constant $v_i$;
the parabola and the other isolated group have distinct constant
$u$ values. High-order Hermite cutoffs in these coordinates
separate the groups. Choose their orders to exceed the largest pole order of the
finitely many basis representatives, measured in the common cover.
Then all cross evaluations have positive order. In a smooth
isolated group the same assertion is ordinary finite jet
interpolation. At the parabola chart, the leading denominator
$-2c$ is a unit; multiplying by a high enough power of
$u-c$ separates the other line without changing that unit.
These constructions give an invertible block leading matrix
of exactly the regular rank. There is no cyclic-seed assumption
and no appeal to a count of unconstructed escape characters.

After a common cover this regular part is trivial, with all
original-marker values of positive order: $t^{-1}$ in the
line and parabola models, $t^{-5/6}$ in the cusp's last model,
$t^{-3/4}$ in the tangent's last model, and the positive powers
of the isolated characters. Lemma~\ref{lem:regularpositiveseed}
excludes a constant original marked period. Differential
isolation and shifting $k$ exclude every nonzero finite
exponential-polynomial in the nonisotropic normal forms.


\subsection{Three concurrent distinct leading lines}
\label{sec:concurrentregular}

Let $F(u,v)$ be a squarefree homogeneous binary cubic,with vertex
the $T$ direction. For the nonisotropic zero-$F_2$ vertex retain the full Gaussian with
$m=q_{TT}\ne0$ and write
\[
 P=T H(u,v)+G(u,v),\qquad\deg H\le1,\quad G_3=F,
 \qquad q=mT^2/2+T J(u,v)+q_*(u,v),
\]
where $J$ is affine linear and $q_*$ includes all lower terms.
For each fixed $t$, the phase gradient is linearly coercive:
its only possible limiting singular direction is the vertex,
where its $T$ derivative has leading term $-mT$. The estimate
is locally uniform on compact parameter sets. Completing $T$
exactly gives
\[
 tP-q=-\frac m2\left(T-\frac{tH-J}{m}\right)^2
                    +tG-q_*+\frac{(tH-J)^2}{2m}.
\]
This change has volume Jacobian one. The binary phase has
squarefree cubic leading form $tF$, whose two leading quadratic
derivatives are a regular sequence of length four. Weighted
leading reduction, or the same Koszul argument used above,
gives full twisted rank four. The Gaussian factor has rank one.
Condition~\eqref{eq:nonisoleading} holds at the sole leading
singular direction. Thus all escapes are excluded from every
pure linear-action block by Proposition~\ref{prop:boundedregularrank}.

If $H$ is a nonzero constant, $P$ has no critical point.
The full regular rank is zero and a constant original marked
period is impossible. This remains true after a constant shift
of $P$.

If $H=0$, integrate the $T$ Gaussian formally. A fixed polynomial
marker becomes a fixed polynomial in $u,v$, using conditional
mean $-J/m$ and covariance $1/m$. The effective Gaussian is
\[
 q_{\rm eff}=q_*-J^2/(2m).
\]
Its Hessian is nondegenerate by the Schur determinant identity
$\det\operatorname{Hess}q=m\det\operatorname{Hess}q_{\rm eff}$.
The squarefree leading binary cubic gives
$\|\nabla G\|\ge c\|(u,v)\|^2$ outside a compact set.
Its finite binary critical set has total Milnor number four.
At a point with Hessian rank at least one,the analytic splitting
lemma gives $\alpha+X^2+Y^{\mu+1}$;Gaussian and power recurrences
produce all $\mu$ characters with prefactors
$t^{-1/2-j/(\mu+1)}$, $1\le j\le\mu$. If the Hessian has rank
zero,the Taylor polynomial is exactly $\alpha+F$ at that point.
Its Milnor number is four,already the full rank. Scaling both
variables by $t^{-1/3}$ gives a fixed binary cubic with rank four
and volume $t^{-2/3}$. Leading-gradient Koszul reduction proves
that rank,and coefficientwise fixed-phase conjugacy constructs
the full frame. All remaining Gaussian terms have positive order.
The binary leading gradient has no projective zero,so there are
no escaping branches under the quadratic perturbation. Polynomial
finite jets separate any distinct critical points. These positive
regular seed orders exclude a constant original marked period;
differential isolation and a constant shift give the stronger
finite exponential-polynomial exclusion. Conditional normalization
is independent of $t$.


Finally let $H$ be nonconstant and use $u=H$ as an affine
coordinate. Then $P=uT+G(u,v)$. If $R(v)=G(0,v)$ is
nonconstant, its critical points give all the bounded groups.
The exact smooth-constraint shear gives their local types
$\alpha+s z+y^{\mu+1}$. This includes the empty critical set
when $R$ is linear. Their full power characters have positive
original-marker orders, so the regular-rank and seed lemmas
exclude a constant period without constructing the escaping
blocks.

If $R=k$ is constant, write $G=k+uB(u,v)$, where $B$ has
degree at most two. The critical curve is
\[
 u=0,\qquad T=-B(0,v).
\]
The quadratic coefficient $c$ of $B(0,v)$ is nonzero.
Otherwise the squarefree cubic $F$ would be divisible by $u^2$.
Put $s=T+B(u,v)$, so exactly $P=k+us$ and the coordinate
Jacobian is one. Scale $u=\rho U$, $s=\rho V$,
$\rho=t^{-1/2}$. The fixed phase is
\[
 UV-q(0,v,-B(0,v)),
\]
with quartic coefficient $-mc^2/2\ne0$ in $v$. Its rank and
critical Milnor number are three. The exact critical equations
$u=\epsilon q_s$, $s=\epsilon q_u$ show that every bounded
branch has bounded scaled coordinates. This is precisely the
complete quartic parabola construction above, now with
Jacobian $\rho^2$. It supplies the entire regular block when
$k=0$, with positive original-marker orders; when $k\ne0$
that block has nonzero exponential part. The polynomial
representatives of $v$ powers are original markers, so its
leading matrix is invertible. The seed lemma applies.

These alternatives exhaust the affine polynomial $H$ and the
restriction $R$. They also exclude nonzero finite
exponential-polynomial markers after differential isolation and
shifting the constant term. We obtain the following conclusion.


\begin{thebibliography}{99}
\bibitem{Pham}
F.~Pham, \emph{La descente des cols par les onglets de Lefschetz,
avec vues sur Gauss--Manin}, \emph{Ast\'erisque} \textbf{130}
(1985), 11--47,
\href{https://numdam.org/item/AST_1985__130__11_0/}{Numdam primary text}.

\bibitem{RoucairolFormal}
C.~Roucairol, \emph{Formal structure of direct image of holonomic
$\mathcal D$-modules of exponential type},
\emph{Manuscripta Mathematica} \textbf{124} (2007), no.~3, 299--318;
\href{https://arxiv.org/abs/math/0604134v2}{arXiv:math/0604134v2}.

\bibitem{SabbahStationary}
C.~Sabbah, \emph{An explicit stationary phase formula for the local
formal Fourier--Laplace transform}, in \emph{Singularities I},
Contemporary Mathematics \textbf{474}, American Mathematical Society,
2008, 309--330;
\href{https://arxiv.org/abs/0706.3570v2}{arXiv:0706.3570v2}.

\end{thebibliography}
\end{document}
